% GENERATED FILE. Edit canonical lessons, then run npm run generate:books.
% canonical-source-hash: fnv1a64:29340b0fa01fa1f2
% canonical-lessons: ZH-C455-xingming, ZH-C455-guoji, ZH-C455-guoqi, ZH-C455-youxiao, ZH-C455-shenqing

\chapter{Forms}
\label{ch:zh-forms}
\hlchaptermodality{455}

\begin{chapteropening}
\textbf{By the end of this chapter:} \emph{I can give my full name and nationality, say a document has expired or is valid, and apply for a visa.}

Complete the last lesson of chapter 455: I can give my full name and nationality, say a document has expired or is valid, and apply for a visa.
\end{chapteropening}

\section[xìngmíng]{xìngmíng — full name}

\label{lesson:ZH-C455-xingming}



\begin{quote}
Before the new one: say the Mandarin for to show someone the way, then the Mandarin for a travel brochure.
\end{quote}

\subsection*{You'll want to know: xìngmíng}
\emph{xìngmíng} — \textquotedblleft{}full name\textquotedblright{}.

\begin{quote}
\emph{qǐng zài zhèr xiě nǐ de xìngmíng} — Please put your full name here.
\end{quote}

Its characters come later; for now it is a word for the ear.

\begin{cousinweb}
\emph{xìngmíng} is built from \emph{xìng} 'a surname' and \zh{名} 'a name', both inherited from Middle Chinese, \emph{sjengH} and \emph{mjieng} (Guangyun, in Baxter's transcription).
\end{cousinweb}

\begin{grammarlens}[title={What you've built}]
One more word for getting around.
\end{grammarlens}

\subsection*{Your turn}
\begin{itemize}
\raggedright
  \item \emph{Say it:} \emph{xìngmíng}
  \item \emph{Say it:} \emph{xìngmíng}, once more
  \item \emph{Say it:} \emph{qǐng zài zhèr xiě nǐ de xìngmíng} — Please put your full name here.
  \item \emph{Recall:} say the Mandarin for to show someone the way, then the Mandarin for a travel brochure, then say \emph{xìngmíng} again
\end{itemize}

\begin{tcolorbox}[breakable,colback=teal!4,colframe=teal!35!black,title={Before you move on}]
What does xìngmíng mean? (\textquotedblleft{}Full name\textquotedblright{}.) Say it once more.
\end{tcolorbox}
\section[guójí]{guójí — citizenship; the country one belongs to}

\label{lesson:ZH-C455-guoji}



\begin{quote}
Before the new one: say the Mandarin for a travel brochure, then the Mandarin for one's full name.
\end{quote}

\subsection*{You'll want to know: guójí}
\emph{guójí} — \textquotedblleft{}citizenship; the country one belongs to\textquotedblright{}.

\begin{quote}
\emph{tā zhù zài Běijīng dànshì tā de guójí shì Yīngguó} — He lives in Beijing, but he's a UK citizen.
\end{quote}

Its characters come later; for now it is a word for the ear.

\begin{cousinweb}
\emph{guójí} is built from \zh{国} 'a country' and \emph{jí} 'a register', both inherited from Middle Chinese, \emph{kwok} and \emph{dzjek} (Guangyun, in Baxter's transcription), which glosses the second as a written roll.
\end{cousinweb}

\begin{grammarlens}[title={What you've built}]
One more word for getting around.
\end{grammarlens}

\subsection*{Your turn}
\begin{itemize}
\raggedright
  \item \emph{Say it:} \emph{guójí}
  \item \emph{Say it:} \emph{guójí}, once more
  \item \emph{Say it:} \emph{tā zhù zài Běijīng dànshì tā de guójí shì Yīngguó} — He lives in Beijing, but he's a UK citizen.
  \item \emph{Recall:} say the Mandarin for a travel brochure, then the Mandarin for one's full name, then say \emph{guójí} again
\end{itemize}

\begin{tcolorbox}[breakable,colback=teal!4,colframe=teal!35!black,title={Before you move on}]
What does guójí mean? (\textquotedblleft{}Citizenship; the country one belongs to\textquotedblright{}.) Say it once more.
\end{tcolorbox}
\section[\texorpdfstring{\zh{过期} (guòqī)}{过期 (guòqī)}]{\zh{过期} (guòqī) — to be out of date; to have expired}

\label{lesson:ZH-C455-guoqi}



\begin{quote}
Before the new one: say the Mandarin for one's full name, then the Mandarin for one's citizenship.
\end{quote}

\subsection*{You'll want to know: \zh{过期}}
\textbf{\zh{过期}} — \emph{guòqī} — \textquotedblleft{}to be out of date; to have expired\textquotedblright{}.

\begin{quote}
\emph{wǒ de qiānzhèng guòqī le} — My visa has expired.
\end{quote}

\begin{cousinweb}
\textbf{\zh{过期}} (\emph{guòqī}) is built from \zh{过} 'to pass' and \zh{期} 'a set time', both inherited from Middle Chinese, \emph{kwaH} and \emph{gi} (Guangyun, in Baxter's transcription), which glosses the second as a time agreed on; today's level-toned qī is irregular beside it.
\end{cousinweb}

\begin{grammarlens}[title={What you've built}]
One more word for getting around.
\end{grammarlens}

\subsection*{Your turn}
\begin{itemize}
\raggedright
  \item \emph{Say it:} \emph{guòqī}
  \item \emph{Say it:} \emph{guòqī}, once more
  \item \emph{Say it:} \emph{wǒ de qiānzhèng guòqī le} — My visa has expired.
  \item \emph{Recall:} say the Mandarin for one's full name, then the Mandarin for one's citizenship, then say \emph{guòqī} again
\end{itemize}

\begin{tcolorbox}[breakable,colback=teal!4,colframe=teal!35!black,title={Before you move on}]
What does \zh{过期} mean? (\textquotedblleft{}To be out of date; to have expired\textquotedblright{}.) Say it once more.
\end{tcolorbox}
\section[yǒuxiào]{yǒuxiào — valid; effective}

\label{lesson:ZH-C455-youxiao}



\begin{quote}
Before the new one: say the Mandarin for one's citizenship, then the Mandarin for to be out of date.
\end{quote}

\subsection*{You'll want to know: yǒuxiào}
\emph{yǒuxiào} — \textquotedblleft{}valid; effective\textquotedblright{}.

\begin{quote}
\emph{zhè zhāng ménpiào jīntiān yǒuxiào} — This ticket is valid today.
\end{quote}

Its characters come later; for now it is a word for the ear.

\begin{cousinweb}
\emph{yǒuxiào} is built from \zh{有} 'to have' and \emph{xiào} 'an effect', both inherited from Middle Chinese, \emph{hjuwX} and \emph{haewH} (Guangyun, in Baxter's transcription), which glosses the second 'to imitate' and 'an effect'.
\end{cousinweb}

\begin{grammarlens}[title={What you've built}]
One more word for getting around.
\end{grammarlens}

\subsection*{Your turn}
\begin{itemize}
\raggedright
  \item \emph{Say it:} \emph{yǒuxiào}
  \item \emph{Say it:} \emph{yǒuxiào}, once more
  \item \emph{Say it:} \emph{zhè zhāng ménpiào jīntiān yǒuxiào} — This ticket is valid today.
  \item \emph{Recall:} say the Mandarin for one's citizenship, then the Mandarin for to be out of date, then say \emph{yǒuxiào} again
\end{itemize}

\begin{tcolorbox}[breakable,colback=teal!4,colframe=teal!35!black,title={Before you move on}]
What does yǒuxiào mean? (\textquotedblleft{}Valid; effective\textquotedblright{}.) Say it once more.
\end{tcolorbox}
\section[shēnqǐng]{shēnqǐng — to apply for}

\label{lesson:ZH-C455-shenqing}



\begin{quote}
Before the new one: say the Mandarin for to be out of date, then the Mandarin for still valid.
\end{quote}

\subsection*{You'll want to know: shēnqǐng}
\emph{shēnqǐng} — \textquotedblleft{}to apply for\textquotedblright{}.

\begin{quote}
\emph{wǒ zuò huǒchē qù Běijīng shēnqǐng qiānzhèng} — I'm taking the train to Beijing to apply for a visa.
\end{quote}

Its characters come later; for now it is a word for the ear.

\begin{cousinweb}
\emph{shēnqǐng} is built from \emph{shēn} 'to state' and \zh{请} 'to ask', both inherited from Middle Chinese, \emph{syin} and \emph{tshjengX} (Guangyun, in Baxter's transcription), which glosses the first 'to stretch' and 'to repeat' and the second 'to beg'.
\end{cousinweb}

\begin{grammarlens}[title={What you've built}]
One more word for getting around.
\end{grammarlens}

\subsection*{Your turn}
\begin{itemize}
\raggedright
  \item \emph{Say it:} \emph{shēnqǐng}
  \item \emph{Say it:} \emph{shēnqǐng}, once more
  \item \emph{Say it:} \emph{wǒ zuò huǒchē qù Běijīng shēnqǐng qiānzhèng} — I'm taking the train to Beijing to apply for a visa.
  \item \emph{Recall:} say the Mandarin for to be out of date, then the Mandarin for still valid, then say \emph{shēnqǐng} again
\end{itemize}

\begin{tcolorbox}[breakable,colback=teal!4,colframe=teal!35!black,title={Before you move on}]
What does shēnqǐng mean? (\textquotedblleft{}To apply for\textquotedblright{}.) Say it once more.
\end{tcolorbox}
